\documentclass[tikz,border=3pt]{standalone}
\usepackage[utf8]{inputenc}
\usepackage[T1]{fontenc}
\usepackage{lmodern}
\renewcommand{\familydefault}{\sfdefault}
\usepackage{sansmath}\sansmath
\usepackage{chemfig}
\usepackage{pgfplots}\pgfplotsset{compat=1.18,/pgf/number format/use comma,/pgf/number format/1000 sep={.}}
\usetikzlibrary{arrows.meta,calc,positioning,decorations.pathmorphing,patterns}
% paleta del sitio
\definecolor{nar}{HTML}{E8680C}\definecolor{narc}{HTML}{FFB35C}\definecolor{hielo}{HTML}{FFF3E8}
\definecolor{tinta}{HTML}{1D1D1F}\definecolor{gris}{HTML}{6E6E73}\definecolor{lila}{HTML}{7C5CF0}
\definecolor{azul}{HTML}{2F6FD6}\definecolor{menta}{HTML}{17A673}\definecolor{coral}{HTML}{E8342A}
\begin{document}
\newif\ifnumeros\numerosfalse
\tikzset{nb/.style={circle,draw=nar,fill=white,text=nar,font=\small\bfseries,inner sep=1.5pt},tx/.style={text=tinta,font=\footnotesize,align=center},
 ovo/.style={draw=lila,fill=lila!25}}
\begin{tikzpicture}[font=\small]
\draw[nar,thick,fill=hielo] (0,0) ellipse (4.2 and 2.6);
% folículos pequeños (primarios)
\foreach \p in {(-3.2,0.5),(-3.0,-0.3),(-2.6,1.2)}{\draw[nar,fill=narc!40] \p circle (0.24); \fill[lila!40] \p circle (0.1);}
% folículo en crecimiento
\draw[nar,fill=narc!45] (-1.4,1.3) circle (0.6);
\fill[azul!10] (-1.25,1.35) circle (0.3);
\draw[ovo] (-1.6,1.15) circle (0.13);
% folículo maduro
\draw[nar,fill=narc!45] (1.0,1.2) circle (1.0);
\fill[azul!12] (1.15,1.25) circle (0.72);
\fill[narc!45] (0.55,0.75) circle (0.32);
\draw[ovo] (0.55,0.8) circle (0.17);
% ovulación
\draw[nar,fill=narc!45] (3.35,0.95) circle (0.55);
\fill[hielo] (3.45,1.05) circle (0.35);
\fill[hielo] (3.6,1.2) -- (4.3,1.9) -- (3.9,0.9) -- cycle;
\fill[narc!50] (4.7,2.2) circle (0.32);
\draw[ovo] (4.7,2.2) circle (0.16);
% cuerpo lúteo
\draw[nar!80!black,fill=narc!90] (2.0,-1.2) circle (0.85);
\foreach \a in {0,45,...,315}{\draw[nar!60!black] (2.0,-1.2) ++(\a:0.25) -- ++(\a:0.45);}
% cuerpo albicans
\draw[gris,fill=gris!15] (-0.9,-1.4) circle (0.5);
% flechas de la secuencia
\draw[gris,dashed,-{Stealth}] (-2.75,0.9) -- (-2.05,1.15);
\draw[gris,dashed,-{Stealth}] (-0.75,1.4) -- (-0.05,1.35);
\draw[gris,dashed,-{Stealth}] (2.05,1.1) -- (2.75,1.0);
\draw[gris,dashed,-{Stealth}] (3.2,0.35) -- (2.65,-0.5);
\draw[gris,dashed,-{Stealth}] (1.1,-1.3) -- (-0.35,-1.4);
\ifnumeros
\draw[gris] (-1.4,1.9) -- (-1.4,2.9) node[nb,anchor=south] {1};
\draw[gris] (1.0,2.2) -- (1.0,2.9) node[nb,anchor=south] {2};
\draw[gris] (4.95,2.4) -- (5.4,2.9) node[nb,anchor=south] {3};
\draw[gris] (2.6,-1.8) -- (3.4,-2.6) node[nb,anchor=north west] {4};
\else
\draw[gris] (-3.0,-0.55) -- (-3.6,-1.7) node[tx,anchor=north] {folículos\\primarios};
\draw[gris] (-1.6,1.85) -- (-2.3,2.9) node[tx,anchor=south] {folículo en\\crecimiento};
\draw[gris] (1.0,2.2) -- (1.0,2.9) node[tx,anchor=south] {folículo maduro};
\draw[gris] (4.95,2.4) -- (5.4,2.9) node[tx,anchor=south] {ovulación: sale el ovocito};
\draw[gris] (2.6,-1.8) -- (3.4,-2.6) node[tx,anchor=north west] {cuerpo lúteo\\(progesterona y estrógeno)};
\draw[gris] (-0.9,-1.9) -- (-1.2,-2.7) node[tx,anchor=north] {cuerpo albicans\\(cicatriz)};
\fi
\end{tikzpicture}
\end{document}
